% Part: computability
% Chapter: computability-theory
% Section: computable-sets

\documentclass[../../../include/open-logic-section]{subfiles}

\begin{document}

\olfileid{cmp}{thy}{cps}
\olsection{قابلِ حساب سیٹ}

قابلِ حساب فنکشناں توں حساب پذیری دا خیال قابلِ حساب
سیٹاں تک ودھایا جا سکدا اے:

\begin{defn}
  من لو $S$ قدرتی عدداں دا اک سیٹ اے۔ تاں $S$ اُدوں تے
  صرف اُدوں \emph{قابلِ حساب} اے جدوں اوہدا خصوصیتی
  فنکشن~$\Char{S}$ قابلِ حساب ہووے۔ ہور لفظاں وچ،
  $S$~اُدوں تے صرف اُدوں قابلِ حساب اے جدوں فنکشن
\[
\Char{S}(x) =
\begin{cases}
1 & \text{جے $x \in S$} \\
0 & \text{ہور صورت وچ}
\end{cases}
\]
قابلِ حساب ہووے۔ ایداں ای تعلق $R(x_0, \dots, x_{k-1})$
اُدوں تے صرف اُدوں قابلِ حساب اے جدوں اوہدا خصوصیتی
فنکشن قابلِ حساب ہووے۔

قابلِ حساب سیٹاں تے تعلقاں نوں \emph{فیصلہ کرن جوگ
(decidable)} وی آکھیا جاندا اے۔
\end{defn}

\begin{explain}
دھیان دیو کہ ہن ساڈے کول حساب پذیری دے کئی خیال نیں:
جزوی فنکشناں لئی، فنکشناں لئی، تے سیٹاں لئی۔ ایہناں
نوں آپس وچ نہ رلاؤ!{} \iftag{TMs}{جزوی فنکشن دا
حساب کرن والی ٹیورنگ مشین اوہناں انپٹ قدراں لئی
فنکشن دا آؤٹ پٹ دیندی اے جنہاں اُتے فنکشن تعریف شدہ
اے؛ سیٹ دا حساب کرن والی ٹیورنگ مشین ایہ فیصلہ کرن
مگروں کہ دِتّی انپٹ قدر سیٹ وچ اے یا نہیں، یا $1$
یا~$0$ واپس دیندی اے۔}{}
\end{explain}

\end{document}
